\chapter*{Introduzione}

Fin dalle rivoluzioni industriali, con i primi macchinari, le fabbriche hanno dovuto gestire i guasti di questi ultimi. Un guasto comporta ritardi alla produzione e relative perdite economiche. Nella storia si è passati dalla manutenzione reattiva, ovvero riparare i macchinari dopo un guasto effettivo, alla manutenzione preventiva programmata, ovvero sostituire le parti più soggette ad usura ad intervalli regolari o a ore-macchina, alla manutenzione predittiva, ovvero sostituire componenti basandosi su misurazioni che possono essere di vibrazione, rumore o temperatura. Questa tesi si concentra sulla manutenzione predittiva.
Le prime modalità di manutenzione predittiva prevedevano l'uso di figure specializzate, chiamate in inglese "machine whisperers" (coloro che sussurrano alle macchine). Questi dipendenti analizzavano i macchinari durante il funzionamento per notare anomalie. Erano molto specializzati, e di conseguenza costosi, anche se per le aziende rappresentavano un costo marginale rispetto ai fermi di produzione e ai costi di riparazione.
A partire dall'Industria 3.0, con i primi macchinari computerizzati si è iniziato ad aggiungere sensori vibrazionali per estrarre informazioni utili a prendere decisioni riguardo alla manutenzione, ma questi sistemi erano ancora dipendenti da persone che analizzassero i dati o da soglie fisse e algoritmi euristici.
L'utilizzo di intelligenze artificiali per l'analisi dei dati viene introdotto con l'Industria 4.0, trasportando i dati nel cloud su server remoti per analizzare i dati storici, tuttavia trasmettere dati grezzi ad alta frequenza verso il cloud satura la banda di rete, introduce latenza e richiede un notevole dispendio energetico e di risorse. Per questo motivo si è iniziato nell'Industria 5.0 a portare l'intelligenza artificiale direttamente nei sensori (\textit{Edge Computing}/\textit{TinyML}).
Secondo la Commissione Europea (rapporto del 2021: "Industry 5.0: Towards a sustainable, human-centric and resilient European industry" \cite{eu_industry50_2021}), l'Industria 5.0 poggia su tre pilastri cardine: Sostenibilità, Resilienza e Uomo al centro. 
Un approccio TinyML aumenta la sostenibilità spostando il calcolo su microcontrollori che costano pochi euro, consumano elettricità nell'ordine dei milliwatt e riducono la trasmissione dei dati tra sensore e cloud, liberando banda. La resilienza aumenta grazie alla capacità del microcontrollore di salvare i dati su una scheda microSD; questo permette all'intero sistema di continuare a monitorare le vibrazioni anche se la connessione cade, per poi caricare in un secondo momento i dati e i metadati verso il cloud. L'uomo rimane comunque centrale al processo decisionale: il servizio non lo sostituisce con una scatola nera, ma funge da strumento di diagnosi precoce, mettendo al primo posto la sicurezza sul lavoro e rilevando in anticipo possibili guasti pericolosi. Inoltre, l'operatore può gestire i dati contrassegnandoli dalla UI se riconosce una condizione specifica, guidando il servizio a creare modelli migliori.
Tuttavia usare un unico modello con un'architettura statica per calcolare le anomalie può risultare controproducente: potrebbe fare fatica a notare anomalie se i dati sono molto eterogenei e potrebbe non essere la migliore architettura per il tipo di dati specifico. Per questo si usano algoritmi NAS (\textit{Neural Architecture Search}), che confrontano varie architetture e dimensioni dei modelli per decidere la più performante. In questa tesi si esplora un'applicazione avanzata di NAS, che unisce tipi diversi di architetture nello stesso gruppo (\textit{ensemble}) con un modello che orchestra diversi sottomodelli; questo permette di scegliere per un macchinario diversi modelli specializzati in stati diversi del macchinario stesso, rendendo il sistema più efficiente ed efficace.
Nella tesi viene esplorato lo sviluppo di un dispositivo capace di raccogliere dati vibrazionali, caricarli nel cloud e ricevere modelli con cui fare l'inferenza sui dati per trovare anomalie; viene inoltre esplorata l'intera infrastruttura cloud a supporto del dispositivo, inclusi la scelta e l'allenamento dei modelli.
